% Generated from locale/tr/content/incompleteness/introduction/definitions.tex; do not hand-edit.


\olfileid[tr]{inc}{int}{def}

\olsection{Tanımlar}

Hilbert'in matematiği biçimselleştirme ve böyle bir biçimselleştirmenin
tutarlı ve tam olduğunu gösterme projesini yürütmek için yapılacak ilk iş
bir dil, mantıksal çerçeve ve aksiyom sistemi seçmek olurdu. Amaçlarımız
bakımından matematiğin birinci derece bir dilde biçimselleştirilebildiğini;
yani birinci derece mantığın bağlaç ve niceleyicileriyle birlikte matematiksel
iddiaları dile getirmemize olanak veren bir !!{constant}s,
!!{function}s ve !!{predicate}s kümesi bulunduğunu varsayalım. Çoğu
kimse böyle bir dilin varlığında uzlaşır: mantıksal olmayan tek simgenin
$\in$ olduğu küme kuramı dili. Böylesine yalın bir dilin bu denli anlatım
gücü taşıdığı ilk bakışta elbette hiç inandırıcı değildir; matematiğin hemen
hemen tamamının bu son derece kıt söz dağarcığında dile getirilebileceğini
kurmak büyük emek gerektirmiştir. Şimdilik işleri yalın tutmak üzere
tartışmamızı aritmetikle, yani matematiğin yalnızca doğal sayılar~$\Nat$ ile
uğraşan bölümüyle sınırlayalım. Aritmetik olgularını dile getirmenin doğal
dili $\Lang L_A$'dır. $\Lang L_A$ tek bir iki yerli !!{predicate}~$<$,
tek bir !!{constant}~$\Obj 0$, bir tek yerli !!{function}~$\prime$ ve iki
iki yerli !!{function}~$+$ ile~$\times$ içerir.

\begin{defn}
Bir !!{sentence}s kümesi~$\Gamma$, gerektirme altında kapalıysa, yani
$\Gamma = \Setabs{!A}{\Gamma \Entails !A}$ ise bir \emph{kuram}dır.
\end{defn}

Kuramları belirtmenin iki kolay yolu vardır. Bunlardan biri, doğru
!!{sentence}s{}i bir !!{structure} içinde toplamak ve kuramı bu küme olarak
vermektir. Örneğin, $\Lang L_A$ için mantıksal olmayan bütün simgelerin
beklendiği gibi yorumlandığı bir !!{structure} ele alın; onun !!{domain}{}i
$\Nat$ olsun.

\begin{defn}\ollabel{def:standard-model}
\emph{Aritmetiğin standart modeli}, aşağıdaki gibi tanımlanan
!!{structure}~$\Struct N$'dir:
\begin{enumerate}
\item $\Domain N = \Nat$
\item $\Assign{\Obj 0}{N} = 0$
\item Her $n \in \Nat$ için $\Assign{\Obj \prime}{N}(n) = n + 1$
\item Her $n, m \in \Nat$ için $\Assign{\Obj +}{N}(n, m) = n + m$
\item Her $n, m \in \Nat$ için $\Assign{\Obj \times}{N}(n, m) = n\cdot m$
\item $\Assign{\Obj <}{N} = \Setabs{\tuple{n, m}}{n \in \Nat, m \in
  \Nat, n < m}$
\end{enumerate}
\end{defn}

`$\times$' ile `$\cdot$' arasındaki ayrıma dikkat edin: $\times$ aritmetik
dilinin bir simgesidir. Onu çarpmayı anımsatsın diye seçtik; ama $\times$
çarpma işlemi değil, iki yerli bir fonksiyon simgesidir (resmî olarak
$\Obj f^2_1$). Buna karşılık $\cdot$, bildiğimiz çarpma fonksiyonunun
\emph{kendisidir}. $n \cdot m$ gibi bir şey gördüğünüzde $n$ ile $m$
sayılarının çarpımını; $x \times y$ gibi bir şey gördüğünüzde ise aritmetik
dilindeki bir terimi kastediyoruz. Standart modelde $\times$ fonksiyon
simgesi, doğal sayılar üzerindeki $\cdot$ fonksiyonu olarak yorumlanır.
Toplama için $+$ simgesini hem aritmetik dilinin fonksiyon simgesi hem de
doğal sayılar üzerindeki toplama fonksiyonu olarak kullanırız. Burada
kastedileni bağlamdan belirlemeniz gerekir.

\begin{defn}
\emph{Doğru aritmetik} kuramı, aritmetiğin standart modelinde sağlanan
!!{sentence}s{}in kümesidir; yani
\[
\Th{TA} = \Setabs{!A}{\Sat{N}{!A}}.
\]
\end{defn}

$\Th{TA}$ bir kuramdır; çünkü $\Th{TA} \Entails !A$ olduğunda $!A$,
$\Th{TA}$'yı sağlayan her !!{structure} içinde sağlanır. $\Sat{N}{\Th{TA}}$
olduğundan $\Sat{N}{!A}$ ve dolayısıyla $!A \in \Th{TA}$ elde ederiz.

Bir $\Gamma$ kuramını belirtmenin öteki yolu, onu bir $\Gamma_0$ tümceler
kümesinin gerektirdiği !!{sentence}s{}in kümesi olarak vermektir. Bu durumda
$\Gamma$, $\Gamma_0$'ın gerektirme altındaki ``kapanışı''dır. Bir kuramı
bu yolla belirtmek ancak $\Gamma_0$ açıkça verilmişse, örneğin
$\Gamma_0$'ın !!{element}s{}ı listelenmişse ilginçtir. En azından
$\Gamma_0$ karar verilebilir olmalıdır; yani !!a{sentence}{}nin
$\Gamma_0$'ın bir elemanı sayılıp sayılmadığını belirleyen hesaplanabilir
bir sınama bulunmalıdır. $\Gamma_0$ içindeki !!{sentence}s{}e $\Gamma$'nın
\emph{aksiyomları}, $\Gamma$'ya da $\Gamma_0$ tarafından
\emph{aksiyomlaştırılmış} deriz.

\begin{defn}
Bir $\Gamma$ kuramı ancak ve ancak
\[
\Gamma = \Setabs{!A}{\Gamma_0 \Entails !A}
\]
ise $\Gamma_0$ tarafından \emph{aksiyomlaştırılmıştır}.
\end{defn}

\begin{defn}
Aşağıdaki tümceler tarafından aksiyomlaştırılan $\Th{Q}$ kuramı
``Robinson $\Th{Q}$'su'' diye bilinir ve çok yalın bir aritmetik kuramıdır.
\begin{align*}
& \lforall[x][\lforall[y][(\eq[x'][y'] \lif \eq[x][y])]] \tag{$!Q_1$}\\
& \lforall[x][\eq/[\Obj 0][x']] \tag{$!Q_2$}\\
& \lforall[x][(\eq[x][\Obj 0] \lor \lexists[y][\eq[x][y']])] \tag{$!Q_3$}\\
& \lforall[x][\eq[(x + \Obj 0)][x]] \tag{$!Q_4$}\\
& \lforall[x][\lforall[y][\eq[(x + y')][(x + y)']]] \tag{$!Q_5$}\\
& \lforall[x][\eq[(x \times \Obj 0)][\Obj 0]] \tag{$!Q_6$}\\
& \lforall[x][\lforall[y][\eq[(x \times y')][((x \times y) + x)]]] \tag{$!Q_7$}\\
& \lforall[x][\lforall[y][(x < y \liff \lexists[z][\eq[(z' + x)][y]])]] \tag{$!Q_8$}
\end{align*}
$\{!Q_1, \dots, !Q_8\}$ !!{sentence}s{}i kümesi $\Th{Q}$'nun
aksiyomlarıdır; dolayısıyla $\Th{Q}$ bunların gerektirdiği bütün
!!{sentence}s{}den oluşur:
\[
\Th{Q} = \Setabs{!A}{\{!Q_1, \dots, !Q_8\} \Entails !A}.
\]
\end{defn}

\begin{defn}
$!A(x)$, $\Lang L_A$ içinde serbest değişkenleri $x$ ve $y_1$, \dots,
$y_n$ olan !!a{formula} olsun. O zaman
\[
\lforall[y_1][\dots\lforall[y_n][((!A(\Obj 0) \land \lforall[x][(!A(x)
\lif !A(x'))]) \lif \lforall[x][!A(x)])]]
\]
biçimindeki her !!{sentence}, \emph{tümevarım şemasının} bir örneğidir.

\emph{Peano aritmetiği}~$\Th{PA}$, $\Th{Q}$'nun aksiyomları ile tümevarım
şemasının bütün örneklerinin birlikte aksiyomlaştırdığı kuramdır.
\end{defn}

\begin{explain}
Tümevarım şemasının her örneği $\Struct{N}$ içinde doğrudur. Bunu görmenin
en kolay yolu, $!A$ !!{formula}{}ünün yalnızca tek bir serbest
!!{variable}{}i~$x$ olduğunu varsaymaktır. Bu durumda $!A(x)$,
$\Struct{N}$ içinde $\Nat$'ın bir $X_{!A}$ altkümesini tanımlar.
$X_{!A}$, $s(x) = n$ iken $\Sat{N}{!A(x)}[s]$ olan bütün $n \in \Nat$
sayılarının kümesidir. Tümevarım şemasının karşılık gelen örneği
\[
((!A(\Obj 0) \land \lforall[x][(!A(x) \lif !A(x'))]) \lif 
  \lforall[x][!A(x)]).
\]
olur. Bunun önceli $\Struct{N}$ içinde doğruysa $0 \in X_{!A}$'dır ve
$n \in X_{!A}$ olduğunda $n+1$ de öyledir. $0 \in X_{!A}$
olduğundan $1 \in X_{!A}$ elde ederiz. $1 \in X_{!A}$ ile $2 \in X_{!A}$
elde ederiz. Böyle devam eder. Dolayısıyla her $n \in \Nat$ için
$n \in X_{!A}$'dır. Bu ise $\lforall[x][!A(x)]$'in $\Struct{N}$ içinde
sağlandığı anlamına gelir.
\end{explain}

Hem $\Th{Q}$ hem de $\Th{PA}$ aksiyomlaştırılmış kuramlardır. Büyük soru,
bunların ne kadar güçlü olduğudur. Örneğin $\Th{PA}$, $\Nat$ hakkındaki
$\Lang L_A$ içinde dile getirilebilen bütün doğruları ispatlayabilir mi?
Daha özel olarak $\Th{PA}$'nın aksiyomları $\Lang L_A$ içinde kurulabilen
bütün soruları karara bağlar mı?

Bunu sormanın başka bir yolu şudur: $\Th{PA} = \Th{TA}$ mıdır? $\Th{TA}$
$\Nat$ hakkındaki bütün doğruları açıkça ispatlar (yani içerir) ve
$\Lang L_A$ içinde kurulabilen bütün soruları karara bağlar; çünkü $!A$,
$\Lang L_A$ içinde !!a{sentence} ise ya $\Sat{N}{!A}$ ya da
$\Sat{N}{\lnot !A}$ olur ve dolayısıyla ya $\Th{TA} \Entails !A$ ya da
$\Th{TA} \Entails \lnot !A$ olur. Böyle bir kurama \emph{!!{complete}}
deyin.

\begin{defn}
Bir $\Gamma$ kuramı \emph{!!{complete}}dır ancak ve ancak kendi dilindeki
her !!{sentence}~$!A$ için ya $\Gamma \Entails !A$ ya da
$\Gamma \Entails \lnot !A$ olur.
\end{defn}

\begin{explain}
Tamlık Teoremi uyarınca $\Gamma \Entails !A$ ancak ve ancak
$\Gamma \Proves !A$ olduğundan, $\Gamma$ ancak ve ancak dilindeki her
!!{sentence}~$!A$ için ya $\Gamma \Proves !A$ ya da
$\Gamma \Proves \lnot !A$ ise tamdır.
\end{explain}

Sormaya yöneldiğimiz başka bir soru şudur: !!a{sentence}{}nin $\Th{TA}$,
$\Th{PA}$, hatta yalnızca $\Th{Q}$ içinde olup olmadığını sınamak için
kullanabileceğimiz hesaplamalı bir yordam var mıdır? Bir kümenin (örneğin
bir !!{sentence}s kümesinin) ne zaman !!{decidable} olduğunu tanımlayarak
bunu daha kesin kılabiliriz.

\begin{defn}
Bir $X$ kümesi, $x$ girdisinde $x \in X$ ise $1$, değilse $0$ döndüren
hesaplamalı bir yordam varsa ancak ve ancak \emph{!!{decidable}} sayılır.
\end{defn}

Öyleyse sorumuz şuna dönüşür: $\Th{TA}$ ($\Th{PA}$, $\Th{Q}$)
!!{decidable} mıdır?

Bu soruların hepsinin yanıtı hayır olacaktır. Bu kuramların hiçbiri karar
verilebilir değildir. Ancak bu olgu bu belirli kuramlara özgü değildir.
Gerçekte belli koşulları sağlayan \emph{her} kuram aynı sonuçlara tabidir.
$\Th{Q}$ ve $\Th{PA}$'nın sağladığı bu koşullardan biri, bunların
!!a{decidable} aksiyomlar kümesiyle aksiyomlaştırılmış olmalarıdır.

\begin{defn}
Bir kuram \emph{!!{axiomatizable}} olarak adlandırılır, eğer
!!a{decidable} aksiyomlar kümesiyle aksiyomlaştırılmışsa.
\end{defn}

\begin{ex}
Sonlu bir !!{sentence}s kümesiyle aksiyomlaştırılan her kuram
!!{axiomatizable} olur; çünkü her sonlu küme !!{decidable} olur. Örneğin
$\Th{Q}$ bu nedenle !!{axiomatizable} bir kuramdır.

$\Th{PA}$ gibi şematik olarak aksiyomlaştırılmış kuramlar da
!!{axiomatizable} olur. Çünkü $!B$'nin $\Th{PA}$'nın aksiyomları arasında
olup olmadığını sınamak, yani $!B$, $\Th{PA}$'nın bir aksiyomuysa
$\Char{X}(!B)=1$, değilse $=0$ olan $\Char{X}$ fonksiyonunu hesaplamak için
şunları yapabiliriz: Önce $!B$'nin $\Th{Q}$ aksiyomlarından biri olup
olmadığını denetleyin. Öyleyse yanıt ``evet''tir ve
$\Char{X}(!B)=1$'dir. Değilse tümevarım şemasının bir örneği olup olmadığını
sınayın. Bu sistemli biçimde yapılabilir; bu durumda belki en kolay şu yolla
yapılabildiği görülür: Tümevarım şemasının her örneği birtakım evrensel
niceleyicilerle, ardından bir koşullu olan alt-!!{formula} ile başlar.
Bu koşullunun ardılı $\lforall[x][!A(x, y_1, \dots, y_n)]$'dir; burada
$x$ ile $y_1$, \dots, $y_n$, $!A$'nın bütün serbest değişkenleridir ve
$!B$'nin başındaki niceleyiciler $y_1$, \dots, $y_n$ değişkenlerini bağlar.
Bu $!A$'yı çıkardıktan ve serbest değişkenlerinin baştaki evrensel
niceleyicilerle ve $\lforall[x]$ ile bağlanan değişkenlere uyduğunu
denetledikten sonra koşullunun öncelinin
\[
!A(\Obj 0, y_1, \dots, y_n) \land \lforall[x][(!A(x, y_1, \dots, y_n)
\lif !A(x', y_1, \dots, y_n))]
\]
ile eşleşip eşleşmediğini denetleriz. Yine, eşleşiyorsa $!B$ tümevarım
şemasının bir örneğidir; eşleşmiyorsa $!B$ böyle bir örnek değildir.
\end{ex}

Bu soruyu ve hangi kuramların tam veya karar verilebilir olduğu yönündeki
daha genel soruyu yanıtlarken aşağıdaki tanımı da ele almak yararlı
olacaktır. Bir $X$ kümesinin !!{enumerable} olduğunu ancak ve ancak boşsa
yahut !!a{surjective} $f \colon \Nat \to X$ fonksiyonu varsa söylemiştik.
Böyle bir fonksiyona $X$'in bir sayımı denir.

\begin{defn}
Bir $X$ kümesine ancak ve ancak boşsa yahut hesaplanabilir bir sayımı varsa
\emph{!!{computably enumerable}} (kısaca !!{c.e.}) denir.
\end{defn}

Kuramların eksiklik teoremlerine tabi olması için !!{axiomatizability}
koşulunun yanında bir başka koşul, hesaplanabilir fonksiyonlar ve !!{decidable}
bağıntılar hakkındaki temel olguları ispatlayacak kadar güçlü olmalarıdır.
``Temel olgular'' ile hesaplanabilir fonksiyonların her argümandaki
değerlerini dile getiren !!{sentence}s{}i kastediyoruz. ``Yeterince güçlü''
ile de söz konusu kuramların bu tümceleri teoremleri arasında saymasını
kastediyoruz. Örneğin tipik bir hesaplanabilir fonksiyonu, toplamayı ele
alın. $+$'nın $2$ ve $3$ argümanlarındaki değeri $5$'tir, yani $2+3=5$.
Aritmetik dilinde $+$'nın $2$ ve $3$ argümanlarındaki değerinin $5$
olduğunu dile getiren bir tümce şudur: $(\num{2}+\num{3})=\num{5}$.
Örneğin $\Th{Q}$ bu tümceyi ispatlar. Daha genel olarak her hesaplanabilir
$f(x_1,x_2)$ fonksiyonu için, $f(n_1,n_2)=m$ olduğunda
$\Th{Q}\Proves !A_f(\num{n_1},\num{n_2},\num{m})$ olacak biçimde
$\Lang{L_A}$ içinde !!a{formula}~$!A_f(x_1,x_2,y)$ bulunmasını isteriz.
Böylece $\Th{Q}$, $f$'nin $n_1$, $n_2$ argümanlarındaki değerinin $m$
olduğunu ispatlar. Gerçekte biraz daha fazlasını, başka hiçbir sayının
$f$'nin $n_1$, $n_2$ argümanlarındaki değeri olmadığını ispatlamasını isteriz.
Aynısı !!{decidable} bağıntılar için de geçerlidir. Bu, aşağıdaki iki tanımda
kesinleştirilir.

\begin{defn}
Bir !!{formula}~$!A(x_1, \dots, x_k, y)$, ancak ve ancak
$f(n_1,\dots,n_k)=m$ olduğunda
\begin{enumerate}
\item $\Gamma \Proves !A(\num{n_1}, \dots, \num{n_k}, \num{m})$ ve
\item $\Gamma \Proves \lforall[y](!A(\num{n_1}, \dots, \num{n_k},
y) \lif y = \num{m})$
\end{enumerate}
oluyorsa $f\colon \Nat^k \to \Nat$ fonksiyonunu $\Gamma$ içinde
\emph{!!{represents} eder}.
\end{defn}

\begin{defn}
Bir !!{formula}~$!A(x_1, \dots, x_k)$ ancak ve ancak
\begin{enumerate}
\item $R(n_1, \dots, n_k)$ olduğunda $\Gamma \Proves !A(\num{n_1},
\dots, \num{n_k})$ ve
\item $R(n_1, \dots, n_k)$ olmadığında $\Gamma \Proves \lnot
!A(\num{n_1}, \dots, \num{n_k})$
\end{enumerate}
oluyorsa $R \subseteq \Nat^k$ bağıntısını \emph{!!{represents} eder}.
\end{defn}

Bir kuram, bütün hesaplanabilir fonksiyonları !!{represents} ediyor ve bütün
!!{decidable} bağıntıları da temsil ediyorsa eksiklik teoremlerinin
uygulanmasına ``yeterince güçlü''dür. $\Th{Q}$ ve onun genişlemeleri bu koşulu sağlar;
ama bunu kurmamız biraz zaman alacaktır. Bu, $\Th{Q}$'nun ne tür şeyleri
ispatlayabildiği hakkında önemsiz olmayan bir olgudur ve $\Th{Q}$'nun bütün
bu olguları ispatlamak için kullanacağımız yalnızca birkaç aksiyomu bulunduğu
için gösterilmesi zordur. Bununla birlikte $\Th{Q}$ çok zayıf bir kuramdır.
Dolayısıyla $\Th{Q}$'nun bütün hesaplanabilir fonksiyonları temsil ettiğini
ispatlamak güç olsa da ilginç kuramların çoğu $\Th{Q}$'dan daha güçlüdür;
yani $\Th{Q}$'nun ispatladığından fazlasını ispatlar. $\Th{Q}$ bir şeyi
ispatlarsa daha güçlü her kuram da onu ispatlar; $\Th{Q}$ bütün
hesaplanabilir fonksiyonları temsil ettiğine göre daha güçlü her kuram da
eder. Bu, birçok ilginç kuramın eksiklik teoremlerinin bu koşulunu karşıladığı
anlamına gelir. Böylece zorlu emeğimiz karşılığını bulacaktır; çünkü
eksiklik teoremlerinin geniş bir kuram sınıfına uygulandığını gösterir.
Elbette ``bütün matematiği'' biçimselleştirmeyi amaçlayan her kuram,
$\Th{Q}$'nun ispatladığı her şeyi ispatlamalıdır; çünkü en azından temel
hesaplamaların sonuçlarını yakalayabilmelidir. Dolayısıyla ``bütün
matematiğin'' kuramı olmaya aday her kuram, eksiklik teoremlerinin
uygulanacağı bir kuram olacaktır.

